\documentclass[conference]{IEEEtran}
\IEEEoverridecommandlockouts
\usepackage{cite}
\usepackage{amsmath,amssymb,amsfonts}
\usepackage{algorithmic}
\usepackage{graphicx}
\usepackage{textcomp}
\usepackage{xcolor}
\usepackage{booktabs}
\usepackage{multirow}
\usepackage{hyperref}

\def\BibTeX{{\rm B\kern-.05em{\sc i\kern-.025em b}\kern-.08em
    T\kern-.1667em\lower.7ex\hbox{E}\kern-.125emX}}

\begin{document}

\title{Degradation Feature Engineering and Classical Machine Learning for Turbofan Engine Remaining Useful Life Prediction\\
\thanks{Advanced Machine Learning (AML) Capstone Project --- Phase 1.}
}

\author{
\IEEEauthorblockN{Sarvesh Choudhary, Arun Jangir, Keshav Chhabra, Mayank Yadav, Rohit Kumar}
\IEEEauthorblockA{\textit{Department of Computer Science and Artificial Intelligence} \\
\textit{Course: Advanced Machine Learning (AML)}\\
GitHub Repository: \url{https://github.com/sarveshDhayal/Capstone_AML}}
}

\maketitle

\begin{abstract}
Accurately estimating the Remaining Useful Life (RUL) of aircraft engines is critical for preventing in-flight failures and optimizing maintenance logistics. In Phase 1 of this investigation, we formulate RUL estimation as a supervised multivariate regression task using the benchmark NASA C-MAPSS dataset (subset FD001). We remove 7 non-informative sensors with zero variance, cap the RUL target at 125 cycles to represent nominal operating life, and extract 71 leakage-free degradation features (moving average, moving standard deviation, trend slope, and deviation from the engine's initial healthy baseline) using past telemetry only. We conduct a rigorous engine-level split (80 training engines, 20 validation engines) to prevent data leakage and evaluate three classical algorithms: Linear Regression (OLS), Decision Tree (CART), and Random Forest. On the official 100 NASA test engines, Random Forest achieves state-of-the-art classical performance with a Test RMSE of 11.44 cycles, MAE of 8.27 cycles, and $R^2$ of 0.918, substantially outperforming the mean baseline (RMSE 41.94 cycles). Engineered features reduce Random Forest validation RMSE by 35.3\% compared to raw sensor telemetry. Feature importance analysis confirms that deviation from the healthy baseline ($\Delta s$) accounts for approximately 85\% of total predictive power, with high-pressure compressor outlet static pressure drift ($s_{11\_dev}$) being the dominant prognostic indicator.
\end{abstract}

\begin{IEEEkeywords}
Prognostics and Health Management (PHM), Remaining Useful Life (RUL), Feature Engineering, Random Forest, Decision Tree, C-MAPSS FD001.
\end{IEEEkeywords}

\section{Introduction}
Gas turbine engines undergo continuous thermomechanical wear during operation. Early detection of subsystem degradation and accurate prediction of Remaining Useful Life (RUL)---the number of operational flight cycles before the engine reaches terminal failure---allows airlines to transition from costly reactive servicing to proactive condition-based maintenance.

The NASA Commercial Modular Aero-Propulsion System Simulation (C-MAPSS) dataset provides realistic run-to-failure multivariate sensor trajectories. However, raw sensor measurements are characterized by high-frequency instrumentation noise and non-stationary operating trends. A fundamental challenge in Prognostics and Health Management (PHM) is transforming raw, noisy multi-channel signals into domain-meaningful degradation indicators that capture wear accumulation without future data leakage.

In this paper (Phase 1), we systematically explore classical machine learning techniques paired with structured degradation feature engineering:
\begin{enumerate}
    \item We formulate a piecewise linear RUL target capped at 125 cycles to address the early-life degradation insensitivity of turbofan sensors.
    \item We engineer 71 leakage-free features capturing temporal smoothing, sensor stability, rate of wear (slope), and drift from each engine's unique healthy baseline ($\Delta s$).
    \item We perform a strict engine-level train/validation split (80/20) and 3-fold GroupKFold cross-validation, guaranteeing zero cross-engine data leakage.
    \item We tune and benchmark Linear Regression, Cost-Complexity Pruned Decision Trees (CART), and Random Forest against the official 100 NASA test engines.
\end{enumerate}

\section{Dataset \& Exploratory Data Analysis}
\subsection{C-MAPSS FD001 Characteristics}
The FD001 dataset models degradation of a high-pressure compressor (HPC) under sea-level static flight conditions:
\begin{itemize}
    \item **Training Set:** 100 engines (20,631 total cycles), all monitored continuously until functional failure.
    \item **Test Set:** 100 engines whose time series are truncated at an arbitrary operational cycle prior to failure.
    \item **Ground Truth ($RUL_{true}$):** Verified cycle counts provided in `RUL\_FD001.txt`.
    \item **Variables (26):** Unit ID, cycle, 3 operational flight settings, and 21 thermodynamic sensors.
\end{itemize}

\subsection{Variance Thresholding \& Flat Sensor Pruning}
Telemetry inspection reveals that 7 sensors exhibit zero variance ($\sigma \approx 0$) across all units and cycles:
\begin{equation}
\{s_1, s_5, s_6, s_{10}, s_{16}, s_{18}, s_{19}\} \rightarrow \text{Pruned}
\end{equation}
Because FD001 simulates a single operating condition, the 3 operational settings are constant and discarded, leaving 14 informative dynamic sensors: $s_2, s_3, s_4, s_7, s_8, s_9, s_{11}, s_{12}, s_{13}, s_{14}, s_{15}, s_{17}, s_{20}, s_{21}$.

\subsection{Engine Lifespans}
Training engine lifespans range from 128 cycles (minimum) to 362 cycles (maximum), with a mean lifespan of $\approx 206$ cycles. Pearson correlation demonstrates that HPC outlet static pressure ($s_{11}$), LPT outlet temperature ($s_4$), and bypass ratio ($s_{15}$) have the strongest linear correlation with degradation ($|r| \in [0.65, 0.70]$).

\section{Methodology \& Feature Engineering}

\subsection{Piecewise Linear Target Formulation}
In physical engines, wear propagation is negligible during the initial 50--100 cycles; the engine behaves as new. A model trained on uncapped linear targets would be forced to predict differing values (e.g., 300 vs 200 cycles) on indistinguishable sensor profiles. We apply a upper cap at $RUL_{cap} = 125$ cycles:
\begin{equation}
y_{i,t} = \min\left(125, \text{cycle}_{i,\max} - t\right)
\end{equation}

\subsection{Engine-Level Data Partitioning}
To prevent optimistic bias, we split the 100 training engines into 80 training units (16,561 records) and 20 validation units (4,070 records) using `GroupShuffleSplit` with seed 42. No engine ID appears in both training and validation sets.

\subsection{Degradation Feature Engineering (71 Features)}
For each of the 14 informative sensors, we compute four features using a 10-cycle backward rolling window ($\tau \in [t-9, t]$) on the same engine:
\begin{enumerate}
    \item **Moving Average ($\mu_s$):** Dampens high-frequency sensor noise:
    \begin{equation}
    \mu_s(t) = \frac{1}{10}\sum_{k=0}^{9} s(t-k)
    \end{equation}
    \item **Moving Standard Deviation ($\sigma_s$):** Detects sensor instability and turbulent flutter:
    \begin{equation}
    \sigma_s(t) = \sqrt{\frac{1}{9}\sum_{k=0}^{9}\left(s(t-k) - \mu_s(t)\right)^2}
    \end{equation}
    \item **Trend Slope ($\beta_s$):** Quantifies rate of wear acceleration via Ordinary Least Squares:
    \begin{equation}
    \beta_s(t) = \frac{\sum_{k=0}^9 (t-k - \bar{t})(s(t-k) - \mu_s(t))}{\sum_{k=0}^9 (t-k - \bar{t})^2}
    \end{equation}
    \item **Baseline Drift ($\Delta s$ or `_dev`):** Difference between current moving average and the engine's initial healthy baseline (mean of cycles $1 \le t \le 10$):
    \begin{equation}
    \Delta s(t) = \mu_s(t) - \frac{1}{10}\sum_{\tau=1}^{10} s(\tau)
    \end{equation}
\end{enumerate}
Together with the current operational cycle and 14 raw sensors, this produces $1 + 14 + (14 \times 4) = 71$ engineered features.

\section{Machine Learning Models \& Tuning}

\subsection{Linear Regression (OLS)}
Acts as the simplest baseline:
\begin{equation}
\hat{y} = w_0 + \sum_{j=1}^{71} w_j x_j
\end{equation}
Features are standardized using `StandardScaler` fitted on the training split.

\subsection{Decision Tree Regressor (CART)}
Captures non-linear piecewise degradation curves. Hyperparameters are tuned via 3-fold `GroupKFold` cross-validation:
\begin{itemize}
    \item `max_depth`: Optimal at depth 7.
    \item `min_samples_leaf`: Optimal at 30 samples.
    \item Post-pruning via cost-complexity pruning ($\alpha = 0.001$).
\end{itemize}

\subsection{Random Forest Regressor}
An ensemble of $B=150$ bootstrap-aggregated (bagged) decision trees. Each tree is trained on a random subset of training engines and random feature subspaces ($\sqrt{p}$), dramatically reducing model variance:
\begin{equation}
\hat{y}_{RF} = \frac{1}{B}\sum_{b=1}^{B} T_b(x)
\end{equation}

\section{Experimental Results and Discussion}

\subsection{Comparative Model Benchmark}
Table~\ref{tab:phase1_results} reports validation and official test set performance across all Phase 1 models.

\begin{table}[htbp]
\caption{Phase 1 Model Performance Benchmark}
\label{tab:phase1_results}
\centering
\begin{tabular}{lcccccc}
\toprule
\multirow{2}{*}{\textbf{Model}} & \multicolumn{3}{c}{\textbf{Validation Set}} & \multicolumn{3}{c}{\textbf{Official NASA Test Set}} \\
\cmidrule(lr){2-4} \cmidrule(lr){5-7}
& \textbf{RMSE} & \textbf{MAE} & \textbf{$R^2$} & \textbf{RMSE} & \textbf{MAE} & \textbf{$R^2$} \\
\midrule
Baseline (Mean) & 41.72 & 36.93 & 0.000 & 41.94 & 34.83 & -0.095 \\
Linear Regression & 15.29 & 11.94 & 0.866 & 15.85 & 12.84 & 0.844 \\
Decision Tree & 14.13 & 9.18 & 0.886 & 14.96 & 10.80 & 0.861 \\
\textbf{Random Forest} & \textbf{10.34} & \textbf{6.77} & \textbf{0.939} & \textbf{11.44} & \textbf{8.27} & \textbf{0.918} \\
\bottomrule
\end{tabular}
\end{table}

Random Forest achieves the superior score on the 100 official test engines: **RMSE 11.44 cycles, MAE 8.27 cycles, and $R^2$ 0.918**. Linear regression struggles because it cannot fit the piecewise flat-then-decay curvature. Single decision trees suffer from high variance, whereas Random Forest's 150-tree bagging eliminates tree instability.

\subsection{Quantitative Effect of Feature Engineering}
To evaluate the impact of the 71 engineered features, we trained identical models using only the 14 raw sensor channels versus with engineered features:
\begin{itemize}
    \item **Linear Regression:** Val RMSE reduced from 17.37 to 15.29 (12.0\% improvement).
    \item **Decision Tree:** Val RMSE reduced from 17.73 to 14.13 (20.3\% improvement).
    \item **Random Forest:** Val RMSE reduced from 15.98 to **10.34** (\textbf{35.3\% improvement}).
\end{itemize}

\subsection{Feature Importance Analysis}
Gini impurity reduction across the 150 trees reveals that:
\begin{itemize}
    \item **Feature Type:** Baseline drift features (`_dev`) account for **84.8\%** of total model importance, followed by moving averages (11.2\%).
    \item **Individual Sensor:** HPC outlet static pressure drift (`s11_dev`) alone contributes **74.1\%** of total importance.
    \item **Secondary Features:** LPT outlet temperature drift (`s4_dev`), bypass ratio drift (`s15_dev`), and core speed drift (`s9_dev`).
\end{itemize}
This demonstrates that measuring an engine's deviation from its \textit{own initial healthy baseline} is substantially more informative than absolute sensor magnitudes.

\subsection{Early vs Late Life Degradation Errors}
Stratifying validation predictions reveals:
\begin{itemize}
    \item **Early Stage ($RUL > 100$):** Random Forest RMSE = 8.01 cycles.
    \item **Late Stage ($RUL \le 100$):** Random Forest RMSE = 12.26 cycles.
\end{itemize}
Late degradation errors are higher due to non-linear physical wear kinetics near the end of life.

\section{Conclusion}
In Phase 1, we proved that domain feature engineering combined with ensemble machine learning provides an effective prognostic framework for turbofan engines. By constructing baseline drift indicators ($\Delta s$) and tuning Random Forests with strict engine-level cross-validation, we achieved an official test RMSE of 11.44 cycles, outperforming raw sensor models by 35.3\%. These results establish the foundational classical benchmark for comparison with Phase 2 deep sequence models.

\begin{thebibliography}{00}
\bibitem{saxena2008} A. Saxena, K. Goebel, D. Simon, and N. Eklund, ``Damage propagation modeling for aircraft engine run-to-failure simulation,'' in \textit{Proc. IEEE PHM}, 2008, pp. 1--9.
\bibitem{breiman2001} L. Breiman, ``Random forests,'' \textit{Machine Learning}, vol. 45, no. 1, pp. 5--32, 2001.
\bibitem{cart1984} L. Breiman, J. Friedman, C. J. Stone, and R. A. Olshen, \textit{Classification and Regression Trees}. CRC press, 1984.
\end{thebibliography}

\end{document}

